% Part: computability
% Chapter: computability-theory
% Section: normal-form

\documentclass[../../../include/open-logic-section]{subfiles}

\begin{document}

\olfileid{cmp}{thy}{nfm}
\olsection{સામાન્ય સ્વરૂપ પ્રમેય}

ધારો કે સંગણનીય વિધેયોની વ્યાખ્યાઓ વર્ણવી શકીએ છીએ અને આપેલા
ઇનપુટ પર તે વિધેયની સંગણનાનો પૂરો ઇતિહાસ દર્શાવતું દેખીતું
વર્ણન સાચું છે કે નહીં તે તપાસી શકીએ છીએ. તો સંગણનાના નિદર્શનથી
સ્વતંત્ર રીતે, કોઈ પણ સંગણનીય વિધેય~$f$નું ઇનપુટ~$x$ પરનું
મૂલ્ય નીચેની રીતે નક્કી કરી શકીએ એવું સ્વાભાવિક લાગે છે:
\begin{enumerate}
  \item સંગણનાના ઇતિહાસનાં શક્ય બધાં વર્ણનોમાં શોધ કરવી.
  \item આપેલો ઇતિહાસ $f(x)$ની સંગણનાનો ઇતિહાસ છે કે નહીં તે તપાસવું.
  \item યોગ્ય ઇતિહાસ મળે તો તેમાંથી $f(x)$નું મૂલ્ય કાઢવું.
\end{enumerate}
ખરેખર આવું થઈ શકે છે તે ક્લીનીના સામાન્ય સ્વરૂપ પ્રમેયનો વિષય છે.

\begin{thm}[ક્લીનીનું સામાન્ય સ્વરૂપ પ્રમેય]
\ollabel{thm:normal-form}
આવો આદિમ પુનરાવર્તી સંબંધ~$T(e,x,s)$ અને આદિમ પુનરાવર્તી
વિધેય~$U(s)$ છે કે દરેક એક-આર્ગ્યુમેન્ટવાળા આંશિક સંગણનીય
વિધેય~$f$ માટે કોઈ~$e$ મળે છે અને દરેક~$x$ માટે
\[
f(x) \simeq U(\umin{s}{T(e,x,s)})
\]
થાય છે.
\sourcecorrection{OLCT-001}{સ્થિર અંગ્રેજી મૂળ ``કોઈ પણ આંશિક
સંગણનીય વિધેય'' કહે છે, પરંતુ પ્રદર્શિત $f(x)$ અને $T(e,x,s)$
એક-આર્ગ્યુમેન્ટવાળું સ્વરૂપ આપે છે. અહીં પ્રમેય એ ક્ષેત્રમાં
સ્પષ્ટ કર્યો છે; વધુ આર્ગ્યુમેન્ટો માટે ટ્યુપલ-કોડ વાપરી શકાય.}
\end{thm}

\begin{proof}[સાબિતીની રૂપરેખા]
સંગણનાના કોઈ પણ નિદર્શન માટે સંગણનીય વિધેય~$f$નું વર્ણન ચુસ્ત રીતે
વ્યાખ્યાયિત કરી શકાય અને તેને પ્રાકૃતિક સંખ્યા~$e$થી સંકેતિત કરી
શકાય. સૂચક~$e$વાળા વિધેયની ઇનપુટ~$x$ પરની સંગણનાની પ્રક્રિયા
નોંધતો ``સંગણના-અનુક્રમ'' પણ ચુસ્ત રીતે વ્યાખ્યાયિત કરી શકાય. એ
અનુક્રમને પણ એક સંખ્યા~$s$થી સંકેતિત કરી શકાય. આ બધું એવી રીતે
કરી શકાય કે
\begin{enumerate}
  \item $T(e,x,s)$ સંબંધ ત્યારે અને માત્ર ત્યારે સાચો હોય જ્યારે
  સંખ્યા~$s$ સૂચક~$e$વાળા વિધેયની ઇનપુટ~$x$ પરની સંગણનાનો
  અનુક્રમ સંકેતિત કરે; અને
  \item $U(s)$ વિધેય $s$થી સંકેતિત સંગણના-અનુક્રમને તે
  સંગણનાના અંતિમ પરિણામ સાથે જોડે.
\end{enumerate}
આ બંને સંગણનીય છે. ખરેખર, $T$ સંબંધ અને $U$ વિધેય આદિમ
પુનરાવર્તી છે.
\end{proof}

\begin{explain}
સામાન્ય સ્વરૂપ પ્રમેયની ચુસ્ત સાબિતી માટે સંગણનાનું એક નિદર્શન
નક્કી કરીને સંગણનીય વિધેયોનાં વર્ણનો તથા સંગણના-અનુક્રમોનું
સંકેતીકરણ વિગતવાર કરવું પડે અને $T$ સંબંધ તથા $U$ વિધેય
આદિમ પુનરાવર્તી છે તે તપાસવું પડે. મોટા ભાગના ઉપયોગોમાં $T$
અને~$U$ સંગણનીય છે તથા $U$ સંપૂર્ણ છે એટલું પૂરતું થાય છે.

પ્રમેયની સાબિતી સૂત્રરૂપે આમ યાદ રાખવી સારી: $\umin{s}{T(e,x,s)}$
સૂચક~$e$વાળા વિધેયની ઇનપુટ~$x$ પરની સંગણનાનો અનુક્રમ શોધે છે,
અને એવો અનુક્રમ મળે તો $U$ તેનું પરિણામ આપે છે.
\end{explain}

જો સંગણનાનું નિદર્શન આંશિક પુનરાવર્તી વિધેયોનું હોય (મૂળે
ક્લીનીએ એ જ વાપર્યું હતું), તો આ પ્રમેય બતાવે છે કે કોઈ પણ
વિધેયની વ્યાખ્યા માટે અનિયત મર્યાદા સુધીની શોધનો એક જ ઉપયોગ,
એટલે એક જ $\umin{y}{f(\vec x,y)}$ ક્રિયા, જરૂરી છે. આ અર્થમાં
દરેક આંશિક પુનરાવર્તી વિધેયને સામાન્ય સ્વરૂપ મળે છે.

$T$ અને $U$ વડે $\cfind{0}$, $\cfind{1}$, $\cfind{2}$, \dots{}ની
પરિગણના વ્યાખ્યાયિત કરી શકાય. હવેથી $T$ અને~$U$ની યોગ્ય પસંદગી
નક્કી થયેલી માનીશું અને
\[
\cfind{e}(x) \simeq U(\umin{s}{T(e,x,s)})
\]
એ સમીકરણને $\cfind{e}$ની \emph{વ્યાખ્યા} ગણીશું.

બીજું ઉપયોગી તથ્ય આ છે:

\begin{thm}
દરેક આંશિક સંગણનીય વિધેયને અનંત ઘણાં સૂચકો હોય છે.
\end{thm}

આ વાત સહજ રીતે પણ સ્પષ્ટ છે. સંગણનીય વિધેયના કોઈ વર્ણનમાંથી
એ જ વિધેય ગણતું બીજું વર્ણન બનાવી શકાય: દાખલા તરીકે, પહેલાં
સંગણનાને અસર ન કરે એવું કંઈક કરવું (જેમ કે $0=0$ છે કે નહીં
તપાસવું અથવા $5$ સુધી ગણવું). બદલાયેલા વર્ણનનો સૂચક મૂળના
સૂચકથી હંમેશાં જુદો હશે. છતાં બંને એક જ વિધેયના સૂચકો છે;
માત્ર તેની સંગણના થોડી જુદી રીતે થાય છે.

\end{document}
